\section{Evaluation}

In this section we present the results evaluation of our algorithm in comparison with the conventional one (as a baseline) and another well-known rational unification algorithm.

We have implemented three unification algorithms with some optimizations\footnote{\url{https://anonymous.4open.science/r/rational-unification-1935}}. These algorithms
were implemented in three variants: ephemeral, ephemeral with path compression and persistent. The latter uses only persistent data structures which makes it possible to utilize them
in a \miniKanren implementation, while others use mutable and thus don't allow backtracking.

To allow fair algorithm comparison all implementations include the same generic optimizations but as many as possible algorithm-specific micro-optimizations.
Especially, we use different data structures for them.

We compare the following algorithms:

\begin{itemize}
    \item \texttt{robinson}~--- the conventional algorithm (Robinson's with triangular substitution, baseline);
    \item \texttt{rational}~--- Martelli-Rossi algorithm (described above);
    \item \texttt{huet}~--- Huet's algorithm for unification of recursive types~\cite{huet1976resolution}\footnote{The actual implementation was studied using \url{https://github.com/SWI-Prolog/swipl/blob/45e83178379f0fd6f918bd294bb361f3db7bf1d7/src/pl-prims.c\#L71}.}.
\end{itemize}

We used \emph{unification traces} for evaluation. Unification trace is a list of pairs of terms, i.e., a sequence of equations. We use three test sets:

\begin{itemize}
    \item \texttt{JGS}~--- collected from the tests from Java Generic Solver implemented in \OCanren~\cite{lozov2023relational} (9 tests, 1448 traces);
    \item \texttt{Lama}~--- collected from the tests from Lama type checker implemented in \OCanren~\cite{domoratskiy2024relational} (66 traces, 1 per test);
    \item \texttt{synth}~--- synthetic tests, generated manually (66 traces, 11 per test).
\end{itemize}

We consider JGS and Lama as unification problem close to a ``real-world'' applications of our algorithm. While JGS includes only unification of finite terms,
several Lama tests require dealing with rational ones.

Synthetic tests are specifically devised for stress testing and parameterized with a certain value $n = 0, 10, \dots, 100$:

\begin{itemize}
    \item \texttt{full} --- equates two full binary trees of depth $n$ ($T_i, U_i \in \V$):
    \begin{equation*}
        \left\{
        \begin{array}{c@{{}={}}ll}
            T_0 & x, & x \in \V, \\
            U_0 & a, & a^0 \in \C, \\
            T_i & f(T_{i-1}, T_{i-1}), & i = 1, 2, \dots, n, \\
            U_i & f(U_{i-1}, U_{i-1}), & i = 1, 2, \dots, n, \\
            T_n & U_n. \\
        \end{array}
        \right.
    \end{equation*}
    \item \texttt{left\_right} --- equates ``leftist'' and ``rightist'' binary trees of depth $n$ ($T_i, U_i \in \V$):
    \begin{equation*}
        \left\{
        \begin{array}{c@{{}={}}ll}
            T_0 & x, \\
            U_0 & x, \\
            T_i & f(T_{i-1}, x), & i = 1, 2, \dots, n, \\
            U_i & f(x, U_{i-1}), & i = 1, 2, \dots, n, \\
            T_n & U_n. \\
        \end{array}
        \right.
    \end{equation*}
    \item $\texttt{loops}_1$ --- equates $n$ loops together:
    \begin{equation*}
        \left\{
        \begin{array}{c@{{}={}}ll}
            x_i & f(x_i), & i = 1, 2, \dots, n, \\
            x_0 & x_i, & i = 1, 2, \dots, n. \\
        \end{array}
        \right.
    \end{equation*}
    \item $\texttt{loops}_2$ --- as the previous but with $x_i = x_0$ instead of $x_0 = x_i$.
    \item $\texttt{loops}_3$ --- as $\texttt{loops}_1$ but with $x_{i-1} = x_i$ instead of $x_0 = x_i$.
    \item $\texttt{loops}_4$ --- as the previous but with $x_i = x_{i-1}$ instead of $x_{i-1} = x_i$.
\end{itemize} 

All synthetic tests except for \texttt{full} require the support for rational unification. Also, \texttt{full} being ran using
the conventional algorithm consumes exponential time. Consequently, we don't use the conventional algorithm for synthetic tests.

% TODO: fix pictures!!!
\begin{figure}[b]
    \centering
    \begin{subfigure}{.48\linewidth}
        \includegraphics[width=\linewidth]{bench/ephemeral/robinson_comparison.png}
        \subcaption{Ephemeral variant}
    \end{subfigure}
    \hfill
    \begin{subfigure}{.48\linewidth}
        \includegraphics[width=\linewidth]{bench/persistent/robinson_comparison.png}
        \subcaption{Persistent variant}
    \end{subfigure}
    \caption{Comparison of the algorithms performance w.r.t. the conventional one}
    \label{fig:baseline}
\end{figure}

The main evaluation results are:

\begin{itemize}
\item Our algorithm performs as fast as the conventional one in the real-world scenarios.
  This is shown by \autoref{fig:baseline}. These plots show the algorithm's run time w.r.t.
  conventional algorithm run time. As we may see, our algorithm is always near to the main
  diagonal.
  
\item While the Huet's algorithm works faster in ephemeral variant, it degrades being implemented
  in persistent settings (this is shown in \autoref{fig:baseline}, too, and in \autoref{fig:Lama}).
  On \autoref{fig:synth-ephemeral} and \autoref{fig:synth-persistent} we may also see that even in
  the ephemeral implementation it works slower than our and moreover its performance depends on the
  ordering of equation sides.
\end{itemize}

\begin{figure}
    \centering
    \begin{subfigure}{.48\linewidth}
        \includegraphics[width=\linewidth]{bench/ephemeral/Lama.png}
        \subcaption{Ephemeral variant}
    \end{subfigure}
    \hfill
    \begin{subfigure}{.48\linewidth}
        \includegraphics[width=\linewidth]{bench/persistent/Lama.png}
        \subcaption{Persistent variant}
    \end{subfigure}
    \caption{Comparison of the algorithms on the \texttt{Lama} traces that require rational unification}
    \label{fig:Lama}
\end{figure}

\begin{figure}
    \centering
    \begin{subfigure}{.48\linewidth}
        \includegraphics[width=\linewidth]{bench/ephemeral/synth_full.png}
        \subcaption{\texttt{full}}
    \end{subfigure}
    \hfill
    \begin{subfigure}{.48\linewidth}
        \includegraphics[width=\linewidth]{bench/ephemeral/synth_left_right.png}
        \subcaption{\texttt{left\_right}}
    \end{subfigure}
    \\
    \begin{subfigure}{.48\linewidth}
        \includegraphics[width=\linewidth]{bench/ephemeral/synth_loops1.png}
        \subcaption{$\texttt{loops}_1$}
    \end{subfigure}
    \hfill
    \begin{subfigure}{.48\linewidth}
        \includegraphics[width=\linewidth]{bench/ephemeral/synth_loops2.png}
        \subcaption{$\texttt{loops}_2$}
    \end{subfigure}
    \\
    \begin{subfigure}{.48\linewidth}
        \includegraphics[width=\linewidth]{bench/ephemeral/synth_loops3.png}
        \subcaption{$\texttt{loops}_3$}
    \end{subfigure}
    \hfill
    \begin{subfigure}{.48\linewidth}
        \includegraphics[width=\linewidth]{bench/ephemeral/synth_loops4.png}
        \subcaption{$\texttt{loops}_4$}
    \end{subfigure}
    \caption{Comparison of the ephemeral rational unification algorithm implementations on the synthetic tests}
    \label{fig:synth-ephemeral}
\end{figure}

\begin{figure}
    \centering
    \begin{subfigure}{.48\linewidth}
        \includegraphics[width=\linewidth]{bench/persistent/synth_full.png}
        \subcaption{\texttt{full}}
    \end{subfigure}
    \hfill
    \begin{subfigure}{.48\linewidth}
        \includegraphics[width=\linewidth]{bench/persistent/synth_left_right.png}
        \subcaption{\texttt{left\_right}}
    \end{subfigure}
    \\
    \begin{subfigure}{.48\linewidth}
        \includegraphics[width=\linewidth]{bench/persistent/synth_loops1.png}
        \subcaption{$\texttt{loops}_1$}
    \end{subfigure}
    \hfill
    \begin{subfigure}{.48\linewidth}
        \includegraphics[width=\linewidth]{bench/persistent/synth_loops2.png}
        \subcaption{$\texttt{loops}_2$}
    \end{subfigure}
    \\
    \begin{subfigure}{.48\linewidth}
        \includegraphics[width=\linewidth]{bench/persistent/synth_loops3.png}
        \subcaption{$\texttt{loops}_3$}
    \end{subfigure}
    \hfill
    \begin{subfigure}{.48\linewidth}
        \includegraphics[width=\linewidth]{bench/persistent/synth_loops4.png}
        \subcaption{$\texttt{loops}_4$}
    \end{subfigure}
    \caption{Comparison of the persistent rational unification algorithm implementations on the synthetic tests}
    \label{fig:synth-persistent}
\end{figure}

We also applied the Kruskal-Wallis test and the Dunn post-hoc test with the Bonferroni correction to determine which
algorithm works faster on which traces. The overall summary is presented as numbers of statistically significant wins 
in \autoref{tab:analyse}. The best results are highlighted using the bold font. Here we see again that the Huet's
algorithm works better in ephemeral implementation but loses its performance in persistent while our stay aligned
with the conventional.

\begin{table}[b]
    \centering
    \begin{tabular}{l||c|c|c}
        \toprule
        \textbf{Algorithm} & \textbf{Ephemeral} & \textbf{Ephemeral with path compression} & \textbf{Persistent} \\
        \midrule
        Conventional & 755 & 915 & \textbf{1982} \\
        Our & 847 & 727 & \textbf{2155} \\
        Huet's & \textbf{2677} & \textbf{2651} & 62 \\
        \bottomrule
    \end{tabular}
    \caption{Statistically significant wins of every algorithm on all traces compared per-algorithm}
    \label{tab:analyse}
\end{table}

Also, the Huet's algorithm requires a more complex data structure representation than the conventional one while our doesn't.
These results show that our algorithm can be used as a drop-in replacement for the conventional one in \miniKanren implementations.

Additionally, following~\cite{domoratskiy2025empirical}, we have studied that path compression technique doesn't
improve overall unification performance. Moreover, it makes our algorithm slightly slower (for about 1\%) while the Huet's slightly
faster. This could be seen in \autoref{tab:analyse-ephemeral}.

\begin{table}[b]
    \centering
    \begin{tabular}{l||c|c|c}
        \toprule
        \textbf{Variant} & \textbf{Conventional} & \textbf{Our} & \textbf{Huet's} \\
        \midrule
        Ephemeral & 597 & \textbf{784} & 494 \\
        Ephemeral with path compression & \textbf{470} & 345 & \textbf{685} \\
        \midrule
        Median ratio & 1.00 & 0.99 & 1.00 \\
        \bottomrule
    \end{tabular}
    \caption{Statistically significant wins of every algorithm on all traces comparing the ephemeral variants}
    \label{tab:analyse-ephemeral}
\end{table}

